\section{Proofs of the main theorems, and worked examples}\label{sec:assembly}

\begin{proof}[Proof of \Cref{thm:lu}]
  Necessity is \Cref{prop:necessity}.  For sufficiency, assume $A$
  satisfies \cref{eq:lu-crit} for all $k$.  By \Cref{lem:reduction-lu},
  $A = M \Pi N$ with $M, N$ invertible triangular and $\Pi$ a matching
  matrix.  By \Cref{cor:transfer-to-normal}, $\Pi$ satisfies
  \cref{eq:lu-crit}, hence by \Cref{lem:comb-lu} the prefix condition
  $\nu_k \le k$ for all $k$.  By \Cref{thm:lu-matching} there are lower
  triangular $L_\Pi$ and upper triangular $U_\Pi$ with
  $L_\Pi U_\Pi = \Pi$.  Then
  \[
    A = M \, (L_\Pi U_\Pi) \, N = (M L_\Pi)(U_\Pi N),
  \]
  and $M L_\Pi$ is lower triangular, $U_\Pi N$ upper triangular.
\end{proof}

\begin{proof}[Proof of \Cref{thm:ldlt}]
  Necessity is \Cref{prop:necessity}.  For sufficiency, assume the
  symmetric $A$ satisfies \cref{eq:ldlt-crit} for all $k$.  By
  \Cref{lem:reduction-sym}, $A = M P M^T$ with $M$ invertible lower
  triangular and $P$ a symmetric matching matrix.  By
  \Cref{cor:transfer-to-normal}, $P$ satisfies \cref{eq:ldlt-crit}, hence
  by \Cref{lem:comb-sym} the prefix condition $\chi_k \le \zeta_k$ for all
  $k$.  By \Cref{thm:ldlt-matching}, $P = S E S^T$ with $S$ lower
  triangular and $E$ diagonal.  Then
  \[
    A = M S \, E \, S^T M^T = (MS) \, E \, (MS)^T,
  \]
  and $MS$ is lower triangular.
\end{proof}

\begin{example}[the exchange matrix, both ways]\label{ex:exchange}
  The matrix $A_2$ of \cref{eq:swap-vs-exchange} is already a normal form.
  As a matching matrix, $\M = \{(2,3), (3,2)\}$ with unit weights; row $1$
  and column $1$ are zero.  The prefix counts are $\nu_1 = 0 \le 1$ and
  $\nu_2 = 2 \le 2$ (both pairs have $\min = 2$), so \Cref{thm:lu-matching}
  applies: the greedy assigns slot $2$ to $(2,3)$ and reroutes $(3,2)$ to
  slot $1$, giving $\ell_2 = e_2$, $u_2 = e_3$, $\ell_1 = e_3$,
  $u_1 = e_2$, which is precisely the LU factorization displayed in
  \cref{eq:exchange-factorizations}.  Index $2$ is a crossing index in the
  sense of \Cref{rem:rerouting}, and slot $1$ --- a zero row \emph{and}
  zero column --- is the resource that resolves it.

  As a symmetric matching matrix, index $1$ is an isolated zero and
  $(2,3)$ is a pair with $\alpha = 1$, $\beta = 0$; the prefix condition
  reads $\chi_2 = 1 \le \zeta_2 = 1$.  There is a single thread
  $1 < 2 < 3$, and the recursion of \Cref{lem:telescope} gives
  ($v_1 = 1$, $\delta_1 = \delta_2 = 1$, $s_1 = -e_2$)
  \[
    L =
    \begin{pmatrix}
      0 & 0 & 0\\
      -1 & -1 & 0\\
      0 & 1 & 1
    \end{pmatrix},
    \qquad
    D = \operatorname{diag}(1, -1, 1),
    \qquad
    L D L^T = A_2 ,
  \]
  the sign-flipped variant of \cref{eq:exchange-factorizations}.  In
  characteristic $2$ the same formulas read $L D L^T$ with
  $D = \operatorname{diag}(1,1,1)$ and $-1 = 1$, and remain valid: no step
  of the construction used signs.
\end{example}

\begin{example}[a genuine crossing]\label{ex:crossing}
  Let $n = 4$ and
  \[
    \Pi =
    \begin{pmatrix}
      0 & 0 & 0 & 0\\
      0 & 0 & 0 & 1\\
      0 & 0 & 0 & 0\\
      0 & 1 & 0 & 0
    \end{pmatrix},
    \qquad
    \M = \{(2,4), (4,2)\},
  \]
  with rows and columns $1, 3$ zero.  Both pairs have $\min = 2$, so
  $\nu_2 = 2 \le 2$; condition (a) of \Cref{lem:comb-lu} reads
  $\chi^R_2 = 1 \le c^0_2 = 1$ and $\chi^R_3 = 1 \le c^0_3 = 2$.  The
  greedy gives slot $2$ to $(2,4)$ and slot $1$ to $(4,2)$:
  \[
    L =
    \begin{pmatrix}
      0 & 0 & 0 & 0\\
      0 & 1 & 0 & 0\\
      0 & 0 & 0 & 0\\
      1 & 0 & 0 & 0
    \end{pmatrix},
    \qquad
    U =
    \begin{pmatrix}
      0 & 1 & 0 & 0\\
      0 & 0 & 0 & 1\\
      0 & 0 & 0 & 0\\
      0 & 0 & 0 & 0
    \end{pmatrix},
    \qquad
    LU = \Pi .
  \]
  Viewed symmetrically, $\Pi$ is the symmetric matching matrix with
  isolated zeros $\{1,3\}$ and the single pair $(2,4)$; the thread
  $1 < 2 < 4$ produces the corresponding $LDL^T$.
\end{example}

\begin{example}[ordering, not rank, is the obstruction]\label{ex:char2}
  Over $GF(2)$, the nonsingular symmetric matrix
  $A = \bigl(\begin{smallmatrix} 0 & 1\\ 1 & 1 \end{smallmatrix}\bigr)$
  fails both criteria at $k = 1$: $\nl(A[1:1,1:1]) = 1$ while the first
  column and row of $A$ are nonzero.  So $A$ has no unpivoted $LU$ or
  $LDL^T$ factorization.  Yet $A$ \emph{is} congruent to the identity by a
  full (non-triangular) congruence:
  with
  $T = \bigl(\begin{smallmatrix} 1 & 1\\ 0 & 1 \end{smallmatrix}\bigr)$,
  $T A T^T = I_2$.  Unpivoted existence is thus a statement about the
  \emph{ordering} of the index set, not about rank or congruence class ---
  consistent with the motivation of preserving a causal or temporal
  ordering discussed in \Cref{sec:intro}.
\end{example}
